The semiaxes are $a=5,b=4$, so the focus is at $F=(3,0)$. Parametrize the ellipse by $Q=(5\cos\theta,4\sin\theta)$. Then
\[
FQ^2=(5\cos\theta-3)^2+16\sin^2\theta
=(5-3\cos\theta)^2.
\]
Thus $FQ=5-3\cos\theta$, whose minimum is $2$, at the endpoint $(5,0)$. A circle centered at $F$ stays inside the ellipse up to this minimum boundary distance, and at radius $2$ it is tangent. The required radius is $\boxed2$.
\begin{asy}
texpreamble("\usepackage{amsmath,amssymb}");
size(240);
defaultpen(fontsize(9pt));
import graph;pair f(real t){return (5cos(t),4sin(t));}draw(graph(f,0,2pi),blue);pair F=(3,0);draw(circle(F,2),red);draw(F--(5,0));dot(F);label("$F$",F,(0,-1));label("$2$",(4,0),(0,1));label("$(5,0)$",(5,0),(1,0));
\end{asy}
